\documentclass[10pt,a4paper]{amsart}
\usepackage[margin=19mm]{geometry}
\usepackage{amsmath,amssymb,mathtools}
\usepackage[scr=rsfs,cal=euler]{mathalfa}
\usepackage{xfrac}
\usepackage[unicode,hidelinks,bookmarks=false]{hyperref}
\newcommand{\ie}{i.e.\ }
\newcommand{\cf}{cf.\ }
\newcommand{\resp}{respectively\ }
\newcommand{\Subsection}{\subsection}
\providecommand{\nobreakdash}{\nobreak\hbox{-}}
\setlength{\emergencystretch}{2em}
\multlinegap0pt
\usepackage{bm}
\usepackage{bbm}
\usepackage{dsfont}
\usepackage{xspace}
\usepackage{cancel}
\usepackage{mathtools}
\usepackage{amsthm}
\makeatletter
\@ifundefined{thm}{\newtheorem{thm}{Theorem}}{}
\@ifundefined{lem}{\newtheorem{lem}[thm]{Lemma}}{}
\makeatother
\title[On the Image of the $p$-adic Logarithm on Annuli of Principa]{On the Image of the $p$-adic Logarithm on Annuli of Principal Units}
\author{Mabud Ali Sarkar}
\date{Source: arXiv:2601.18187v1 (CC BY 4.0)}
\begin{document}
\maketitle

\noindent\emph{Source and licence.} On the Image of the $p$-adic Logarithm on Annuli of Principal Units. Version arXiv:2601.18187v1 (CC BY 4.0), distributed under the Creative Commons Attribution 4.0 International licence, as recorded in the arXiv licence metadata for this exact version and shown on the abstract page https://arxiv.org/abs/2601.18187v1. Full rights notice: licenses/arxiv-2601.18187-CC-BY-4.0.txt.\par\smallskip

\noindent\textit{Editorial scope.} Counted unit: the Lem whose source label is \texttt{A1} in the article (source0.tex lines 238--253, source label \texttt{A1}), delivered together with the article's own complete original proof of it (source0.tex lines 255--326, one \texttt{proof} environment of 72 lines). No mathematics is written, repaired or re-derived: statement and proof are the article's own text line for line. Required context retained uncounted: t2.1 (lines 129--131). Every definition, display and label of those lines is retained unchanged. Omitted from this package and not redistributed: 1--128, 132--237, 254--254, 327--328. Nothing inside a retained excerpt is omitted or altered: every retained body is delivered exactly as the article's lines are written, so no omission receipt of any kind accompanies this package. Citations: the retained excerpts carry no citation command at all, so no bibliography entry of the article is transcribed and no reference is redistributed. Reference adaptation: none was needed and none was made; every reference inside the delivered unit resolves inside it, so no unresolved reference or citation is produced. Printed slips kept literal: no printed word, symbol, hypothesis or number of the counted statement or of its proof was altered. Layout and class: the article's own class is \texttt{amsart} with packages including pgfplots, tikz, color, amsmath, amsfonts, amssymb, amsthm, graphicx, soul, appendix, lipsum, epstopdf, pdflscape, csquotes; the wrapper is the lane's pinned \texttt{amsart} editorial wrapper at 10pt on A4, which restates the theorem-like environments the retained text needs and adds the article's own macro definitions that the retained text needs (none), with extra wrapper packages (bm, bbm, dsfont, xspace, cancel, mathtools). The delivered numbering is the unit's own. Assets: no retained span contains a figure environment, a graphics-inclusion command, a PDF-page inclusion, a table or a TikZ picture; no image file is copied into this package, the package declares no asset dependency and no image file is redistributed with it. Licence: version arXiv:2601.18187v1 of this article is distributed under the Creative Commons Attribution 4.0 International licence, as recorded in the arXiv licence metadata for this exact version and shown on the abstract page https://arxiv.org/abs/2601.18187v1. A full rights notice is delivered at licenses/arxiv-2601.18187-CC-BY-4.0.txt. Characters: every retained excerpt is pure ASCII, so the delivered engine, XeLaTeX with its default Latin Modern text font, reports no missing character.
\begin{thm}\label{t2.1}
Let $p \geq 3 $ be a prime and consider the cyclotomic extension $K=\mathbb{Q}_p(\zeta_p)$ of $\mathbb{Q}_p$, where $\zeta_p$ is a primitive $p$th root of unity, and $\mathfrak{m}_K$ be the maximal ideal. Then the image of $(1+\mathfrak{m}_K) \setminus (1+\mathfrak{m}_K^2)$ under $p$-adic logarithm is $\mathfrak{m}_K^2$.
\end{thm}

\begin{lem} \label{A1}
Let $K=\mathbb{Q}_p(\zeta_p)$ with uniformizer $\pi$ satisfying $\pi^{p-1}=-p$, and let
\[
y=\sum_{n\ge 2} y_n \pi^n \in \mathfrak m_K^2 .
\]
There exists an element
\[
a=1+\sum_{n\ge 1} a_n \pi^n \in (1+\mathfrak m_K)\setminus(1+\mathfrak m_K^2),
\qquad a_n\in\{0,1,\dots,p-1\},
\]
such that
\[
\log_p(a)=y .
\]
Moreover, once $a_1\neq 0$ is fixed, the coefficients $a_n$ can be determined inductively.
\end{lem}

\begin{proof}
Let us write
\[
x=\sum_{n\ge 1} a_n \pi^n \in \mathfrak m_K ,
\qquad a=1+x .
\]
Since $v(x)>0$, the $p$-adic logarithm admits the convergent expansion
\[
\log_p(1+x)=x-\frac{x^2}{2}+\frac{x^3}{3}-\cdots .
\]


Modulo $\pi^3$, we computes
\[
\log_p(1+a_1\pi+a_2\pi^2)
\equiv a_1\pi+\left(a_2-\frac{a_1^2}{2}\right)\pi^2
\pmod{\pi^3}.
\]
As shown in the proof of Theorem~\ref{t2.1}, the coefficient of $\pi$ lies in
$\mathfrak m_K^2$, hence the condition
\[
\log_p(a)\equiv y_2\pi^2 \pmod{\pi^3}
\]
reduces to
\[
a_2-\frac{a_1^2}{2}\equiv y_2 \pmod{p}.
\]
Given $a_1 \neq 0$, this congruence admits solutions, and for each such choice the
coefficient $a_2$ is uniquely determined modulo $p$.

Assume that for some $N\ge 2$ the coefficients
$a_1,\dots,a_N$ have been chosen so that
\[
\log_p\!\left(1+\sum_{n=1}^N a_n\pi^n\right)
\equiv \sum_{n=2}^N y_n\pi^n
\pmod{\pi^{N+1}}.
\]
Set
\[
a^{(N+1)}=1+\sum_{n=1}^N a_n\pi^n+a_{N+1}\pi^{N+1}.
\]
Then
\[
a^{(N+1)}=a^{(N)}\left(1+u\right),
\qquad
u=\frac{a_{N+1}\pi^{N+1}}{a^{(N)}},
\]
with $v(u)=N+1$. Using the additivity of $\log_p$ on sufficiently small elements,
we obtain
\[
\log_p(a^{(N+1)})
\equiv \log_p(a^{(N)}) + a_{N+1}\pi^{N+1}
\pmod{\pi^{N+2}}.
\]
Hence
\[
\log_p(a^{(N+1)})
\equiv \sum_{n=2}^N y_n\pi^n
+ \left(c_{N+1}+a_{N+1}\right)\pi^{N+1}
\pmod{\pi^{N+2}},
\]
where $c_{N+1}\in\mathcal O_K$ is the error term, depends only on $a_1,\dots,a_N$.
The congruence
\[
c_{N+1}+a_{N+1}\equiv y_{N+1}\pmod{p}
\]
has a unique solution for $a_{N+1}\in\{0,1,\dots,p-1\}$, which completes the
inductive step.


The resulting series converges in $\mathcal O_K$, satisfies $a_1\neq 0$, and by induction $\log_p(a)=y$.
\end{proof}

\end{document}
